\subsection{A parity improvement for prescribed-point rational covers}
\label{sec:rank-one-augmentation}

\paragraph{Finding A26 (the rank-one augmentation criterion).}
Let $C/\mathbf F_Q$ be a smooth proper geometrically connected
genus-two curve in characteristic two, let $P\in C(\mathbf F_Q)$,
and suppose $N=\#C(\mathbf F_Q)$ is even.
Let $r$ be an odd prime with $\operatorname{ord}_r(2)=r-1$, and set
\[
 G=J_C(\mathbf F_Q),\qquad V=G/rG.
\tag{RA.1}
\]
If $|V|>N$, some surjection $G\twoheadrightarrow\mathbf Z/r$
gives a constant, geometrically connected, $P$-split cyclic
\'etale Lang cover whose relative class number is odd.
If $r\mid|G|$ and $r>N$, every such surjection has this property.
In both statements the prescribed point $P$ is arbitrary.
No lower bound on $N$ in terms of the kernel size is required.

Every genus-two curve of $2$-rank one satisfies the stated
parity hypothesis over its finite field of definition and
over every finite extension.  Thus both assertions apply
to every such curve whenever their remaining numerical
hypotheses hold.

\begin{proof}
We first prove the parity assertion.  The hyperelliptic
degree-two map in characteristic two is an Artin--Schreier
cover.  If its reduced geometric poles have odd orders
$d_1,\ldots,d_t$, its genus and $2$-rank are
\[
       g=\frac{\sum_{i=1}^t(d_i+1)-2}{2},
       \qquad f_C=t-1;
\]
see \cite[Lemma~2.6, p.~712]{PriesZhu2012}.
For $g=2$ and $f_C=1$, there are two poles, with
orders $1$ and $3$.  Frobenius preserves the ramification
orders, so it fixes both branch points individually.
Each has one geometric point above it, which is therefore
rational.  These are exactly the two fixed points of the
hyperelliptic involution.  Every other rational point is
paired with a distinct rational point by that involution.
Consequently $N$ is even.  The same argument applies over
every extension; in particular the two branch points remain
rational.

Now assume only the evenness of $N$.
In the odd-order group algebra $\mathbf F_2[V]$, put
\[
 \Theta_V=[0]+\sum_{R\in C(\mathbf F_Q)}[[R-P]\bmod rG].
\tag{RA.2}
\]
Its augmentation is $N+1=1$ in $\mathbf F_2$.
Recovered A1 B.2 gives the following exact criterion:
for a surjection $a:G\to\mathbf Z/r$, the associated
$P$-split Lang cover has odd relative class number if and
only if every nontrivial character of $\mathbf Z/r$ has
nonzero evaluation on $a(\Theta_V)$
\cite[B.2]{A1run}.
The Lang construction there also gives geometric
connectedness and the rational point above the prescribed
$P$.

Suppose all surjections to $\mathbf Z/r$ fail the criterion.
The irreducibility of $\Phi_r$ over $\mathbf F_2$, which
follows from $\operatorname{ord}_r(2)=r-1$, means that
one vanishing nontrivial character evaluation for such a
quotient forces all its nontrivial character evaluations
to vanish.  Every nontrivial character of $V$ factors
through one of these quotients.  Thus all nontrivial
Fourier coefficients of $\Theta_V$ vanish.
Fourier inversion is valid because $|V|$ is odd, and gives
\[
        \Theta_V=c\sum_{b\in V}[b],\qquad
        c=\frac{N+1}{|V|}=1\quad\text{in }\mathbf F_2.
\tag{RA.3}
\]
But $|V|>N$ leaves at least one empty Abel-image coset.
The zero coset is occupied by $P$, so an empty coset is
nonzero.  Its coefficient in (RA.2) is zero, contradicting
(RA.3).  Therefore some quotient succeeds.

For the assertion about every surjection, fix one
$a:G\to\mathbf Z/r$ and suppose it fails.
Apply the same argument to
$\Theta_a=[0]+\sum_R[a([R-P])]$ in
$\mathbf F_2[\mathbf Z/r]$.
Irreducibility and the augmentation force
$\Theta_a=\sum_{b\in\mathbf Z/r}[b]$.
Since $r>N$, an empty nonzero coset again contradicts
that equality.  As $a$ was arbitrary, every quotient
succeeds.  A surjection exists when $r\mid|G|$ by the
structure theorem for finite abelian groups.
\end{proof}

\paragraph{Relation to the recovered row statements.}
The Fourier identity (RA.3) is already the displayed
identity in recovered successor A.35, and the empty-coset
argument also occurs in A.10 \cite[A.10 and A.35]{ASrun}.
Those statements retain an additional pair-count
inequality; the present even-$N$ specialization obtains
an immediate contradiction from augmentation and does
not require that inequality.
Recovered A1 B.1 then gives the prescribed-point
separation implication for the resulting odd cyclic cover.
This is a rational relative-class-number condition over
the displayed field, not geometric primitive $2$-rank
zero; it need not survive an arbitrary constant extension.
The existence of a prime satisfying (RA.1) and $|V|>N$
is an additional requirement, not a consequence of
Hooley's positive density statement alone.
